\subsection{Moser isotopy and Darboux's theorem}

\blemm[title={Moser isotopy, compact}, name={mosercompactisotopy}]

    Let $M$ be a compact manifold, and $\omega_t$ a smooth family of symplectic forms on $M$ which has an exact derivative:
    $$ \drv t\omega_t = \d\sigma_t,\qquad \sigma_t\in\Omega^1(M) . $$
    Then there exists an isotopy $\psi_t\in\Diff(M)$ such that $\psi_t^*\omega_t=\omega_0$.

\elemm

\bproof

    Instead of considering isotopies, we consider the time-dependent vector fields $X_t$ which are associated to them.
    Because $M$ is compact, this is a one-tone-one correspondence.
    So we have
    $$ \drv t\psi_t = X_t\circ\psi_t,\qquad \psi_0 = \id . $$
    Now differentiate $\psi_t^*\omega_t$, we want it to be zero:
    $$ \drv t\psi_t^*\omega_t = \psi_t^*\parens{\drv t\omega_t + \L_{X_t}\omega_t} = \psi_t^*\parens{\d\sigma_t + \d\iota(X_t)\omega_t} . $$
    So it is sufficient to require that
    $$ \sigma_t + \iota(X_t)\omega_t = 0 , $$
    so $X_t=\hat\omega_t^{-1}\sigma_t$ gives the time-dependent vector field.
    Taking its flow gives the desired isotopy.
    \qed

\eproof

We want to extend this result to non-compact manifolds.
As-is we cannot do this, but the following lemma is sufficient.

\blemm[title={Moser isotopy}, name={moserisotopy}]

    Let $M$ be a manifold and $Q\subseteq M$ a compact submanifold.
    Suppose that $\omega_0,\omega_1\in\Omega^2(M)$ are closed $2$-forms such that $\omega_0$ and $\omega_1$ agree on $Q$ and are nondegenerate at each point of $q$.
    Then there exist open neighborhoods $U_0,U_1$ of $Q$ and a diffeomorphism $\psi\colon U_0\to U_1$ such that
    $$ \psi|_Q = \id,\qquad \psi^*\omega_1 = \omega_0 . $$

\elemm

\bproof

    Note that if $\omega$ is nondegenerate at a point $p\in M$, then it is nondegenerate in a neighborhood of $p$.
    Indeed, relative to any local frame of $p$ $\omega$ has nonsingular representation at $p$, and taking this representation is smooth, so the result follows
    because $\gl(n)$ is open in $\Mat_n$.
    Thus $\omega_0,\omega_1$ are both nondegenerate in an open neighborhood $U_0$ of $Q$.

    We claim the existence of a $1$-form $\sigma\in\Omega^1(U_0)$ such that
    $$ \sigma|_Q = 0,\qquad \d\sigma = \omega_1-\omega_0 . $$
    Then define
    $$ \omega_t = \omega_0 + t(\omega_1-\omega_0) = \omega_0 + t\d\sigma . $$
    This, as a convex combination of $\omega_0,\omega_1$, is nondegenerate on $U_0$ as well.
    Applying Moser's argument above to $Q$, we get diffeomorphisms $\psi_t\colon U_0\to U_t$ satisying $\psi_t^*\omega_t=\omega_0$; take $\psi=\psi_1$ for the lemma.
    In more detail, define $X_t=\hat\omega_t^{-1}\sigma$ in $U_0$, then this time-dependent vector field has global flow in a neighborhood of $Q$ as it is compact.
    Shrink $U_0$ to be this neighborhood, so the $\psi_t$ are diffeomorphisms from $U_0$ to a neighborhood of $Q$.

    Thus all that remains is to find a $\sigma\in\Omega^1(U_0)$ which restricts to $0$ on $Q$ and whose exterior derivative is $\omega_1-\omega_0$.
    Endow $M$ with any Riemannian metric, then $Q$ has a tubular neighborhood.
    We explain what this means.
    Let $\exp\colon\E\subseteq TM\to M$ be the exponential map, where $\exp(p,tv)$ is the geodesic starting at $p$ with initial velocity $v$.
    Recall that $\E$ is an open subset of $TM$ whose fibers are star-shaped about $0$.
    Let $TQ^\perp\subseteq TM|_Q$ be the orthogonal subbundle of $TQ\subseteq TM|_Q$, then there is an open subset $W\subseteq TQ^\perp$ of the form
    $$ W = \set{(q,v)}[q\in Q,\ v\in T_qQ^\perp,\ \norm v<\epsilon] $$
    such that $\exp$ forms a diffeomorphism from $W$ to $\exp(W)$ which is also open.
    This $\exp(W)$ is called a {\emphcolor tubular neighborhood} of $Q$; its existence is nontrivial.

    Define $U_0$ to be the tubular neighborhood $\exp(W)$.
    Then define for $0\leq t\leq1$,
    $$ \phi_t\colon U_0\to U_0,\qquad \phi_t(\exp(q,v)) = \exp(q,tv) , $$
    which is a diffeomorphism for $t>0$.
    And $\phi_0(\exp(q,v))=\exp(q,0)=q$ so that $\phi_0(U_0)\subseteq Q$, $\phi_1=\id$, and $\phi_t|_Q=\id$.
    The last claim is because $\exp(q,v)\in Q$ iff $v=0$ because $\exp$ is a diffeomorphism from $W$ (because $\exp(q,v)=\exp(\exp(q,v),0)$).

    Let $\tau=\omega_1-\omega_0$, then these facts mean
    $$ \phi_0^*\tau = 0,\qquad \phi_1^*\tau = \tau . $$
    For $t>0$ define
    $$ X_t = (\partial_t\phi_t)\circ\phi_t^{-1} . $$
    Then we have
    $$ \drv t\phi_t^*\tau = \phi_t^*\L_{X_t}\tau = \d(\phi_t^*(\iota(X_t)\tau)) , $$
    so define $\sigma_t=\phi_t^*(\iota(X_t)\tau)$.
    Then for $w\in T_xU_0$ we have
    $$ \sigma_t(x;w) = \tau(\phi_t(x);\ X_t(\phi_t(x)),\d\phi_t(x)(w)) = \tau(\phi_t(x);\ \partial_t\phi_t(x),\d\phi_t(x)(w)) . $$
    This is then smooth at $t=0$, and for $x\in Q$ vanishes (because $\phi_t(x)=x$ and $\tau(x;\bul)=0$ as $\omega_0,\omega_1$ agree on $Q$).
    Thus
    $$ \omega_1 - \omega_0 = \tau = \phi_1^*\tau - \phi_0^*\tau = \int_0^1\drv t\parens{\phi_t^*\tau}\d t = \int_0^1\d\sigma_t\d t , $$
    so define
    $$ \sigma = \int_0^1\sigma_t\d t $$
    then we get
    $$ \omega_1 - \omega_0 = \d\sigma . $$
    And on $Q$, the $\sigma_t$ vanish and thus so too does $\sigma$, so $\sigma$ satisfies the conditions.
    \qed

\eproof

\bthrm[title={Darboux's}, name={darboux}]

    Every symplectic form $\omega$ on $M$ is locally symplectomorphic to the standard form $\omega_0$ on $\bR^{2n}$.

\ethrm

\bproof

    Choose a point $q\in M$, then $\omega$ has a symplectic basis on $T_qM$.
    Let $(x_i,y_i)$ be a local coordinate chart on $q\in U$ such that $(\partial_ix_i,\partial_iy_i)$ is a symplectic basis of $\omega_0$ on $T_qM$.
    Define the symplectic form on $U$,
    $$ \tilde\omega = \sum_i\d x_i\wedge\d y_i . $$
    By \refmath{moserisotopy}, by taking $Q=\set q$, there is a diffeomorphism $\psi\colon U_0\to U_1$, with $q\in U_0,U_1\subseteq U$ open, such that $\omega=\psi^*\tilde\omega$.
    So $(U_0,\omega)$ is symplectomorphic to $(U_1,\tilde\omega)$, which is clearly symplectomorphic to the standard Euclidean symplectic form $\omega_0$.
    \qed

\eproof

\bnote

    In Riemannian geometry, not every metric is locally isomorphic to the Euclidean metric.
    The failure of this is measured by the curvature of the metric.
    In contrast to this, Darboux's theorem shows that all symplectic manifolds are locally isomorphic.

    Thus there are no meaningful local invariants for symplectic manifolds, and so symplectic geometry focuses on global invariants.

\enote

\bcoro

    Let $(M,\omega)$ be a symplectic manifold of dimension $2n$.
    Then there exists an open cover $\set{U_\alpha}_\alpha$ of $M$ with charts $\alpha\colon U_\alpha\to\alpha(U_\alpha)\subseteq\bR^{2n}$
    such that $\alpha^*\omega_0=\omega$.
    The atlas has symplectic transition matrices:
    $$ \d(\beta\circ\alpha^{-1})(x) \in \Sp(2n),\qquad x\in\alpha(U_\alpha\cap U_\beta) . $$
    Such charts are called {\emphcolor Darboux charts}.

\ecoro

\bthrm[title={Moser stability}, name={moserstability}]

    Let $M$ be a closed manifold, and suppose $\omega_t$ is a smooth family of cohomologous symplectic forms on $M$.
    Then there is a smooth family of diffeomorphisms $\psi_t$ of $M$ such that
    $$ \psi_0 = \id,\qquad \psi_t^*\omega_t = \omega_0 . $$

\ethrm

\bproof

    We apply \refmath[lemma]{mosercompactisotopy}, so we need to find a smooth family of $1$-forms $\sigma_t$ such that
    $$ \d\sigma_t = \tau_t = \drv t\omega_t . $$
    Because $\omega_t$ are cohomologous, $\omega_t-\omega$ are exact and thus so too is $\tau_t$.
    Thus there do exist $1$-forms $\sigma_t$ such that $\d\sigma_t=\tau_t$.
    The only issue is ensuring $\sigma_t$ is smooth, and this can be done by induction on the number of sets in a good cover of $M$.
    \TODO{This?}
    \qed

\eproof

\bdefn

    Two symplectic forms $\omega_0,\omega_1$ on a manifold $M$ are {\emphcolor isotopic} if they can be connected by a smooth family
    of cohomologous symplectic forms $\omega_t$.

    They are said to be {\emphcolor strongly isotopic} if there is an isotopy $\psi_t$ of $M$ such that $\psi_1^*\omega_0=\omega_1$.

\edefn

In general being strongly isotopic implies being isotopic.
This is because we can define $\omega_t=\psi_t^*\omega_0$.
And as $\psi_\bul$ is a homotopy, by the homotopy-invariance of de Rham cohomology, $\psi_t^*$ induces the same map on homotopy for all $t$.
\refmath{moserstability} says that the two notions coincide when $M$ is closed.

